% META: {"id": "1NSI-ADGK-CO-001", "chapitre": "1NSI-ALGO-DICHO-GLOUTON-KNN", "type_objet": "corrige", "exercice_ref": "1NSI-ADGK-EX-001", "capacites": ["P-ALGO-04"], "sources_inspiration": [], "status": "generated"}
% BEGIN-VERIFY
% def recherche_dichotomique_trace(tableau_trie, valeur):
%     borne_inf, borne_sup = 0, len(tableau_trie) - 1
%     etapes = []
%     while borne_inf <= borne_sup:
%         milieu = (borne_inf + borne_sup) // 2
%         etapes.append((borne_inf, borne_sup, milieu, tableau_trie[milieu]))
%         if tableau_trie[milieu] == valeur:
%             return milieu, etapes
%         elif tableau_trie[milieu] < valeur:
%             borne_inf = milieu + 1
%         else:
%             borne_sup = milieu - 1
%     return -1, etapes
% t = [2, 5, 8, 12, 16, 23, 38, 45, 56, 72]
% idx, etapes = recherche_dichotomique_trace(t, 23)
% assert idx == 5
% assert etapes == [(0, 9, 4, 16), (5, 9, 7, 45), (5, 6, 5, 23)]
% variants = [sup - inf for (inf, sup, _, _) in etapes]
% assert variants == [9, 4, 1]
% for k in range(1, len(variants)):
%     assert variants[k] < variants[k - 1]
% END-VERIFY
\begin{corrige}{1NSI-ADGK-EX-001}
\textbf{1.}
\begin{center}
\begin{tabular}{cccc}
  \lstinline{borne_inf} & \lstinline{borne_sup} & \lstinline{milieu} & \lstinline{tableau_trie[milieu]} \\\hline
  0 & 9 & 4 & 16 \\
  5 & 9 & 7 & 45 \\
  5 & 6 & 5 & 23 \\
\end{tabular}
\end{center}

\textbf{2.} $3$ itérations sont nécessaires. L'indice renvoyé est $5$.

\textbf{3.} $V$ vaut successivement $9-0=9$, $9-5=4$, $6-5=1$ : la suite $9,4,1$ est bien
strictement décroissante.
\end{corrige}
